% Part: first-order-logic
% Chapter: axiomatic-deduction
% Section: proving-things-quant

\documentclass[../../../include/open-logic-section]{subfiles}

\begin{document}

\olfileid{fol}{axd}{prq}

\olsection{પરિમાણકોવાળી \usetoken{P}{derivation}}

\begin{ex}
$(\lforall[x][!A(x)] \land \lforall[y][!B(y)]) \lif
\lforall[x][(!A(x) \land !B(x))]$ની !!a{derivation} આપીએ.

પહેલાં ધ્યાન આપો કે
\begin{align*}
  (\lforall[x][!A(x)] \land \lforall[y][!B(y)]) & \lif \lforall[x][!A(x)]\\
  \intertext{એ \olref[prp]{ax:land1}નો દૃષ્ટાંત છે, અને}
  \lforall[x][!A(x)] & \lif !A(a) \\
  \intertext{એ \olref[qua]{ax:q1}નો દૃષ્ટાંત છે. તેથી \olref[pro]{prop:chain} વડે જાણીએ છીએ કે}
  (\lforall[x][!A(x)] \land \lforall[y][!B(y)]) & \lif !A(a) \\
  \intertext{!!{derivable} છે. એ જ રીતે, કારણ કે}
  (\lforall[x][!A(x)] \land \lforall[y][!B(y)]) & \lif \lforall[y][!B(y)] \qquad\text{અને}\\
    \lforall[y][!B(y)] & \lif !B(a)\\
    \intertext{અનુક્રમે \olref[prp]{ax:land2} અને \olref[qua]{ax:q1}ના દૃષ્ટાંતો છે,}
    (\lforall[x][!A(x)] \land \lforall[y][!B(y)]) & \lif !B(a)\\
    \intertext{એ \olref[pro]{prop:chain} વડે નિષ્પન્ન કરી શકાય છે. \olref[prp]{ax:land3}નો યોગ્ય દૃષ્ટાંત અને~\MP{}ના બે ઉપયોગ લઈને દેખાય છે કે}
    (\lforall[x][!A(x)] \land \lforall[y][!B(y)]) & \lif (!A(a) \land !B(a))\\
    \intertext{નિષ્પન્ન કરી શકાય છે. હવે \QR{} લાગુ કરતાં મળે છે}
(\lforall[x][!A(x)] \land \lforall[y][!B(y)]) & \lif \lforall[x][(!A(x) \land !B(x))].
\end{align*}
\end{ex}

\end{document}
